% Generated draft: requires mathematical and source review.
\documentclass[11pt]{article}
\usepackage[margin=25mm]{geometry}
\usepackage{fontspec}
\setmainfont{Latin Modern Roman}
\usepackage{amsmath,amssymb,mathtools,unicode-math}
\setmathfont{Latin Modern Math}
\usepackage{enumitem,needspace,xurl,hyperref}
\setlength{\parindent}{0pt}
\setlength{\parskip}{5pt}
\setlength{\emergencystretch}{3em}
\newcommand{\N}{\mathbb{N}}
\newcommand{\Z}{\mathbb{Z}}
\newcommand{\Q}{\mathbb{Q}}
\newcommand{\R}{\mathbb{R}}
\newcommand{\C}{\mathbb{C}}
\newcommand{\F}{\mathbb{F}}
\newcommand{\A}{\mathbb{A}}
\newcommand{\PP}{\mathbb{P}}
\newcommand{\E}{\mathbb{E}}
\newcommand{\eps}{\varepsilon}
\newcommand{\dd}{\,\mathrm{d}}
\begin{document}
\section*{UnsolvedMath ID: 700569}
\textbf{Source problem code:} catalogue entry
\subsection{Definitions and mathematical statement}
For a topological group \(G\), let
\[
\mathcal{C}(G)=\{H\leq G:H\text{ is closed in }G\},
\]
ordered by inclusion. This is a complete lattice: for any family \((H_i)_{i\in I}\) of closed subgroups,
\[
\bigwedge_{i\in I}H_i=\bigcap_{i\in I}H_i,
\qquad
\bigvee_{i\in I}H_i=\overline{\left\langle\bigcup_{i\in I}H_i\right\rangle}^{\,G}.
\]
Here \(\langle S\rangle\) denotes the subgroup generated by \(S\), and the bar denotes closure in \(G\). The empty meet is \(G\), and the empty join is \(\overline{\{1_G\}}^{\,G}\), which equals \(\{1_G\}\) when \(G\) is Hausdorff.

The problem is to characterize, by necessary and sufficient conditions on a lattice \(L\), when there exist a locally compact topological group \(G\) and a lattice isomorphism
\[
L\cong\mathcal{C}(G).
\]
A lattice isomorphism here means a bijection preserving and reflecting order. This formulation reads the requested characterization as one of abstract lattices; the source specifies no additional topology on the lattice. The source also does not explicitly state whether locally compact groups must be Hausdorff, so that convention remains unresolved and must be fixed for a definitive characterization.

The source reports that the discrete case was characterized by B. V. Yakovlev in the cited 1974 publication. In that case every subgroup is closed, so the lattice in question is the ordinary subgroup lattice. This historical assertion is supplied source information, not an independently verified assessment.
\subsection{Short English statement}
Characterize the abstract lattices that occur as lattices of closed subgroups of locally compact groups, ordered by inclusion. The source reports an existing characterization for discrete groups. Its Hausdorff convention is unspecified.
\subsection{Sources}
\begin{itemize}
\item Supplied local UnsolvedMath registry record: permanent ID \(700569\), exact source code \texttt{catalogue entry}, title \emph{Kourovka Notebook Problem 9.36}. The registry labels the problem \texttt{open}; this is supplied metadata, not independently verified mathematical status.
\item \emph{The Kourovka Notebook (Unsolved Problems in Group Theory)}, problem 9.36, proposed by Yu. N. Mukhin, 9th issue (1984). The supplied background also identifies the 21st issue edition (2026), \texttt{arXiv:1401.0300v48}. The collection description names E. I. Khukhro and V. D. Mazurov as editors.
\item B. V. Yakovlev, \emph{Algebra and Logic}, volume 13, no. 6 (1974), pp. 400--412. The source cites this publication for the discrete case; it supplies no article title.
\end{itemize}
\end{document}
